\documentclass[11pt,a4paper]{article}
\usepackage[utf8]{inputenc}
\usepackage[T1]{fontenc}
\usepackage{geometry}
\geometry{a4paper, margin=0.75in}
\usepackage[colorlinks=true, urlcolor=blue, linkcolor=blue]{hyperref}
\usepackage{titlesec}
\usepackage{enumitem}
\usepackage{xcolor}

% Set font to Helvetica or standard sans-serif
\renewcommand{\familydefault}{\sfdefault}

% Section formatting
\titleformat{\section}{\large\bfseries}{}{0em}{}[\titlerule]
\titlespacing*{\section}{0pt}{14pt}{6pt}

\pagestyle{empty} % Suppress page numbers for a clean look

\begin{document}

% Header
\begin{center}
    {\LARGE \textbf{ Nikita Kazeev }} \\
    \vspace{4pt}
    \small
    \href{mailto:kazeevn@gmail.com}{ kazeevn@gmail.com } $\cdot$ 
    \href{ https://kazeevn.github.io }{Website} $\cdot$ 
    \href{ https://github.com/kazeevn }{GitHub} $\cdot$ 
    \href{ https://linkedin.com/in/nikita-kazeev }{LinkedIn} $\cdot$ 
    \href{ https://scholar.google.com/citations?user=vamy2okAAAAJ }{Google Scholar} $\cdot$ 
    \href{ https://orcid.org/0000-0002-5699-7634 }{ORCID}
\end{center}

\vspace{6pt}

% Summary/Title
\begin{center}
    \textit{ Research Scientist at the intersection of AI and Physics }
\end{center}

% Work Experience Section
\section*{Work Experience}

\noindent \textbf{ Perfomax } \hfill \textbf{ May 2026--present } \\
\textit{ AI Advisor }
\begin{itemize}[leftmargin=*, noitemsep, topsep=3pt]

  \item \textbf{ AI Strategy \& Platform Architecture }: 
  \begin{itemize}[leftmargin=*, label=$\cdot$, noitemsep, topsep=1pt]
  
    \item Advised on AI/ML strategy, product implementation, and technical design for the platform.
  
  \end{itemize}

\end{itemize}
\vspace{8pt}


\noindent \textbf{ National University of Singapore } \hfill \textbf{ 2022--present } \\
\textit{ Postdoc under Kostya Novoselov }
\begin{itemize}[leftmargin=*, noitemsep, topsep=3pt]

  \item \textbf{ Leadership }: 
  \begin{itemize}[leftmargin=*, label=$\cdot$, noitemsep, topsep=1pt]
  
    \item Speaker search and selection for the 500-person AI4X 2025 conference.
  
    \item Main organizer of the ICLR 2025 workshop on multiscale machine learning.
  
  \end{itemize}

  \item \textbf{ Transformer for generating symmetric crystals }: 
  \begin{itemize}[leftmargin=*, label=$\cdot$, noitemsep, topsep=1pt]
  
    \item Created a generative model for material design based on symmetry inductive bias.
  
    \item Implemented the model in PyTorch.
  
    \item Achieved the best generation quality and diversity among space-group-conditioned models.
  
    \item Led a team of 6.
  
  \end{itemize}

  \item \textbf{ ML for predicting properties of defects in 2D crystals }: 
  \begin{itemize}[leftmargin=*, label=$\cdot$, noitemsep, topsep=1pt]
  
    \item Created a sparse representation of crystals with defects, achieving 4x energy prediction error reduction.
  
    \item Developed code for reproducible parallel running of machine learning experiments.
  
    \item Led a team of 3.
  
  \end{itemize}

\end{itemize}
\vspace{8pt}


\noindent \textbf{ Constructor Group } \hfill \textbf{ February 2021--present } \\
\textit{ Machine Learning Consultant }
\begin{itemize}[leftmargin=*, noitemsep, topsep=3pt]

  \item \textbf{ AI/ML Consulting }: 
  \begin{itemize}[leftmargin=*, label=$\cdot$, noitemsep, topsep=1pt]
  
    \item Consulted on machine learning algorithms, deep learning applications, and architectural designs for scientific and industry-led projects.
  
  \end{itemize}

\end{itemize}
\vspace{8pt}


\noindent \textbf{ CERN [Yandex $\rightarrow$ HSE University] } \hfill \textbf{ 2014--2022 } \\
\textit{ Researcher } \\ \small \href{https://breakthroughprize.org/Laureates/1/L3995}{2025 Breakthrough Prize in Fundamental Physics} as a member of LHCb.
\begin{itemize}[leftmargin=*, noitemsep, topsep=3pt]

  \item \textbf{ Muon identification at the LHCb experiment at CERN }: 
  \begin{itemize}[leftmargin=*, label=$\cdot$, noitemsep, topsep=1pt]
  
    \item Solved a classification problem over tabular data with noisy labels under timing constraints.
  
    \item Developed a gradient boosting model that reduced the false positive rate by 30\% in the critical low-track-momentum region.
  
    \item Integrated the model into LHCb production, mostly in C++.
  
    \item Packaged the work as a data science competition problem for IDAO-2019.
  
  \end{itemize}

  \item \textbf{ CatBoost aka fighting biases with dynamic boosting }: 
  \begin{itemize}[leftmargin=*, label=$\cdot$, noitemsep, topsep=1pt]
  
    \item Set up a distributed system for running experiments for the team with Bash and Python.
  
    \item Studied gradient boosting improvements with experiments on toy data.
  
  \end{itemize}

\end{itemize}
\vspace{8pt}



% Education Section
\section*{Education}

\noindent \textbf{ Higher School of Economics (HSE University) } \hfill \textbf{ 2016--2020 } \\
PhD in Computer Science, supervisor \href{https://www.linkedin.com/profile/view?id=2422149}{Andrey Ustyuzhanin}; \href{https://repository.cern/records/034yx-nt191}{Machine Learning for particle identification in the LHCb detector}.
\vspace{6pt}


\noindent \textbf{ Sapienza --- Universita di Roma } \hfill \textbf{ 2016--2020 } \\
PhD in Physics (double degree with HSE), supervisor \href{https://it.linkedin.com/in/barbara-sciascia-234935a}{Barbara Sciascia}.
\vspace{6pt}


\noindent \textbf{ Product Management course at Yandex } \hfill \textbf{ Spring 2018 } \\
Focused coursework in product strategy and execution.
\vspace{6pt}


\noindent \textbf{ Yandex School of Data Analysis } \hfill \textbf{ 2013--2015 } \\
Master's level CS course covering algorithms, machine learning, deep learning, and distributed systems.
\vspace{6pt}


\noindent \textbf{ Moscow Institute of Physics and Technology } \hfill \textbf{  } \\
\textbf{MS in Physics} (2014--2016): Optimisation of data processing of the LHCb experiment.

\textbf{BS in Physics} (2010--2014): Study of the quantum states of the electrons in nonideal plasma and selected molecules using wave packet molecular dynamics with packet splitting.
\vspace{6pt}



% Mentorship Section

\section*{Mentorship}
\begin{itemize}[leftmargin=*, noitemsep]

  \item Mentored 9 students (\href{https://www.hse.ru/edu/vkr/219430466}{1}, \href{https://www.hse.ru/edu/vkr/296294066}{2}, \href{https://www.hse.ru/edu/vkr/219524130}{3}, \href{https://www.hse.ru/en/edu/vkr/470561417}{4}, \href{https://www.hse.ru/en/edu/vkr/470935224}{5}, \href{https://www.hse.ru/ma/datasci/students/diplomas/631565328}{6}, \href{https://www.hse.ru/ma/datasci/students/diplomas/631565422}{7}, \href{https://www.hse.ru/edu/vkr/624890096}{8}, \href{https://www.hse.ru/ba/ami/students/diplomas/834351454}{9}), 3 interns, and student workshop projects (\href{https://sochisirius.ru/news/3148}{1}, \href{https://www.hse.ru/edu/vkr/296294066}{2}0).

  \item Student council member from 2013 to 2016, leading several dormitory-scale IT projects.

  \item Co-PI of a \$3.4m AI Singapore grant on multiscale machine learning.

\end{itemize}


% Teaching Section

\section*{Teaching \& Outreach}
\begin{itemize}[leftmargin=*, noitemsep]

  \item Pitch research to the general public through major newspaper coverage (\href{https://www.kommersant.ru/doc/7832307}{1}, \href{https://www.kommersant.ru/doc/6122817}{2}), blog posts (\href{https://ifim.nus.edu.sg/tailor-made-2d-materials-became-closer-with-the-new-machine-learning-algorithm-developed-at-nus/}{1}, \href{https://communities.springernature.com/posts/sparse-representation-for-machine-learning-the-properties-of-defects-in-2d-materials}{2}), a \href{https://elementy.ru/events/435506/IV\_Nauchnye\_boi\_NIU\_VShE\_na\_Dne\_Vyshki}{science slam}, interviews (\href{https://web.archive.org/web/20210518043835/https://academy.yandex.ru/posts/beresh-i-razbiraeshsya}{1}, \href{https://habr.com/en/company/yandex/blog/422761/}{2}), and a public \href{https://www.youtube.com/watch?v=zLFZ\_z-7-JI&feature=youtu.be&t=7779}{talk}.

  \item Taught at the Machine Learning in High Energy Physics Summer School in \href{https://indico.cern.ch/event/439520/}{2015}, \href{https://indico.cern.ch/event/497368/}{2016}, \href{https://indico.cern.ch/event/613571/}{2017}, \href{https://indico.cern.ch/event/687473/}{2018}, \href{https://indico.cern.ch/event/768915/}{2019}, \href{https://indico.cern.ch/event/838377/}{2020}, \href{https://indico.cern.ch/event/1025052/}{2021}, and \href{https://indico.cern.ch/event/1229514/}{2023}.

  \item Lecturer in the \href{https://www.coursera.org/specializations/aml}{Coursera Advanced Machine Learning specialisation}, which is \href{https://www.hse.ru/en/news/edu/259378493.html}{award-winning}.

  \item Taught and developed an introductory data analysis \href{https://we.hse.ru/en/}{course} at HSE.

  \item Delivered more than 20 \href{https://docs.google.com/document/d/1IOlBse4q0YgEoL\_RCJTvJfbFhFaD\_YMAIJJ3W77B\_X0/edit?usp=sharing}{conference talks}.

\end{itemize}


% Service Section

\section*{Service}
\begin{itemize}[leftmargin=*, noitemsep]

  \item Reviewer for RSC Advances, Machine Learning: Science and Technology, and AI for Accelerated Materials Design workshops at NeurIPS and ICLR (2023--2025).

  \item Peer Staff Supporter at NUS, serving as a first-line support for mental wellbeing.

\end{itemize}


% Skills Section

\section*{Skills}
\noindent

\textbf{ ML \& AI } (Structured, non-text, domain-specific data) $\cdot$ 

\textbf{ Python } (Research coding) $\cdot$ 

\textbf{ PyTorch } (Deep learning) $\cdot$ 

\textbf{ Linux } (Administration) $\cdot$ 

\textbf{ Research Writing \& Speaking }  $\cdot$ 

\textbf{ Leadership \& Mentorship } 



% Publications Section
\section*{Selected Publications}
\begin{enumerate}[leftmargin=*, noitemsep]

  \item \textbf{Kazeev, Nikita}, and Phan Bui Nhat Huyen. \href{https://openreview.net/forum?id=J3Vqv2fFse}{Position: Align AI to Our Aspirations, Not Our Flaws}. Pluralistic Alignment Workshop at ICML 2026, June 2026.

  \item Artem Dembitskiy, Shuya Yamazaki, Artem Maevskiy, \textbf{Nikita Kazeev}, Roman A. Eremin, Semen Budennyy, Kedar Hippalgaonkar, Antonio Helio Castro Neto, and Andrey E. Ustyuzhanin. \href{https://openreview.net/forum?id=xFI6K4ZQDp}{Generative Pipeline for Discovering Solid-State Battery Materials with Universal Atomistic Potentials}. ICML 2026 AI for Science Workshop, May 2026.

  \item Jose Recatala-Gomez, \textbf{Nikita Kazeev}, et al. \href{https://arxiv.org/abs/2604.14082}{Generative design of inorganic materials}. arXiv preprint arXiv:2604.14082, April 2026.

  \item Alexandre Duval, Siddharth Betala, Samuel P. Gleason, Andy Xu, Georgia Channing, Daniel Levy, Ali Ramlaoui, Clémentine Fourrier, Chaitanya K. Joshi, \textbf{Nikita Kazeev}, Sékou-Oumar Kaba, Félix Therrien, Alex Hernández-García, Rocío Mercado, N M Anoop Krishnan. \href{https://arxiv.org/abs/2512.04562}{LeMat-GenBench: Bridging the gap between crystal generation and materials discovery}. AI for Accelerated Materials Design - NeurIPS 2025 / arXiv preprint arXiv:2512.04562, October 2025.

  \item Stephen G. Dale, \textbf{Nikita Kazeev}, et al. \href{https://arxiv.org/abs/2511.20976}{AI4X Roadmap: Artificial Intelligence for the advancement of scientific pursuit and its future directions}. arXiv preprint arXiv:2511.20976, November 2025.

  \item \textbf{Kazeev, Nikita}, et al. \href{https://arxiv.org/abs/2503.02407}{WyckoffTransformer: Generation of Symmetric Crystals}. ICML 2025, May 2025.

  \item \textbf{Kazeev, N.}, Al-Maeeni, A.R., Romanov, I. et al. \href{https://doi.org/10.1038/s41524-023-01062-z}{Sparse representation for machine learning the properties of defects in 2D materials}. npj Comput Mater 9, 113 (2023).

\end{enumerate}

\end{document}